% ============================================================================
% BAB II - ARSITEKTUR DAN TOPOLOGI SISTEM
% ============================================================================
\chapter{ARSITEKTUR DAN TOPOLOGI SISTEM}

\section{Diagram Topologi Jaringan}

Arsitektur sistem terdiri dari 5 (lima) \textit{node} utama yang
berjalan di atas jaringan Docker \textit{bridge} dengan subnet
\texttt{172.25.0.0/24}. Diagram pada Gambar~\ref{fig:topologi}
menggambarkan hubungan antar \textit{node} beserta alamat IP, port
layanan, dan arah komunikasi data yang telah dirancang. Setiap
\textit{node} memiliki peran spesifik --- mulai dari \textit{client}
browser, aplikasi Laravel, middleware Pgpool-II, hingga tiga
\textit{backend} PostgreSQL (satu Master dan dua Replica). Jaringan
Docker \textit{bridge} dipilih karena menyediakan isolasi jaringan
internal yang aman dengan resolusi nama \textit{host} otomatis antar
\textit{container}. Konfigurasi IP statis diterapkan untuk memastikan
setiap \textit{node} selalu dapat dijangkau pada alamat yang konsisten
meskipun terjadi restart \textit{container}.

\begin{figure}[H]
\centering
\resizebox{\textwidth}{!}{%
\begin{tikzpicture}[
    node distance=0.8cm,
    box/.style={rectangle, draw=black!70, fill=white, rounded corners=2pt,
                minimum width=3.2cm, minimum height=0.9cm, align=center,
                font=\footnotesize\sffamily, blur shadow={shadow blur steps=4}},
    titlebox/.style={rectangle, draw=black!40, fill=black!5, rounded corners=2pt,
                     minimum width=3.2cm, minimum height=0.65cm, align=center,
                     font=\footnotesize\bfseries\sffamily},
    dbnode/.style={rectangle, draw=blue!60, fill=blue!3, rounded corners=2pt,
                   minimum width=2.8cm, minimum height=0.85cm, align=center,
                   font=\footnotesize\ttfamily},
    arrow/.style={-{Stealth[scale=0.9]}, line width=0.7pt, draw=black!60},
    writearrow/.style={-{Stealth[scale=0.9]}, line width=1.0pt, draw=red!70, thick},
    readarrow/.style={-{Stealth[scale=0.9]}, line width=1.0pt, draw=green!50!black, thick},
    replarrow/.style={-{Stealth[scale=0.9]}, line width=1.0pt, draw=blue!60, thick, dashed},
    labeltext/.style={font=\footnotesize\sffamily, fill=white, inner sep=1pt},
]

% --- Client ---
\node[box, minimum width=2cm] (client) at (0, 0) {Browser\\\scriptsize (Client)};

% --- Laravel App ---
\node[box, right=1.0cm of client] (app) {\textbf{Laravel App}\\bdt-app:80\\{\scriptsize 172.25.0.30}};

% --- Pgpool-II ---
\node[box, below=1.8cm of app, minimum height=1.2cm] (pgpool) {\textbf{Pgpool-II}\\pgpool:5432 (DB) / :9898 (PCP)\\{\scriptsize 172.25.0.20}};

% --- Master ---
\node[dbnode, below left=2.0cm and 0.3cm of pgpool] (master) {\underline{PostgreSQL}\\\underline{Master}\\pg-master:5432\\{\scriptsize 172.25.0.10}};

% --- Replica 1 ---
\node[dbnode, below=2.0cm of pgpool] (repl1) {\underline{PostgreSQL}\\\underline{Replica 1}\\pg-replica1:5432\\{\scriptsize 172.25.0.11}};

% --- Replica 2 ---
\node[dbnode, below right=2.0cm and 0.3cm of pgpool] (repl2) {\underline{PostgreSQL}\\\underline{Replica 2}\\pg-replica2:5432\\{\scriptsize 172.25.0.12}};

% --- Arrows: Client to App ---
\draw[arrow] (client.east) -- (app.west) node[labeltext, midway, above] {\scriptsize HTTP :8080};

% --- Arrow: App to Pgpool ---
\draw[arrow] (app.south) -- ++(0, -0.15) -| (pgpool.north) node[labeltext, pos=0.6, above right] {\scriptsize DB Connection};

% --- Arrow: Pgpool to Master (WRITE) ---
\draw[writearrow] (pgpool.south west) -- ++(-0.3, -0.2) |- (master.north) node[labeltext, midway, left] {\scriptsize WRITE};

% --- Arrow: Pgpool to Replica 1 (READ) ---
\draw[readarrow] (pgpool.south) -- (repl1.north) node[labeltext, midway, right] {\scriptsize READ};

% --- Arrow: Pgpool to Replica 2 (READ) ---
\draw[readarrow] (pgpool.south east) -- ++(0.3, -0.2) |- (repl2.north) node[labeltext, midway, right] {\scriptsize READ};

% --- Arrows: Master to Replicas (Streaming Replication) ---
\draw[replarrow] (master.east) -- ++(0.6, 0) |- (repl1.west) node[labeltext, midway, above, rotate=90] {\scriptsize Streaming Replication};
\draw[replarrow] (repl1.east) -- (repl2.west) node[labeltext, midway, above] {\scriptsize Streaming Replication};

% --- Host boundary ---
\node[draw=black!30, dashed, rounded corners=5pt, fit=(client)(app)(pgpool)(master)(repl1)(repl2),
      inner sep=0.5cm, label={[font=\footnotesize\sffamily]north east:{\scriptsize Host Machine (Docker)}}] (host) {};

\end{tikzpicture}
}
\caption{Diagram Topologi Jaringan Sistem \textit{Read-Write Splitting}}
\label{fig:topologi}
\end{figure}

\subsection{Alamat IP dan Port Layanan}

Tabel~\ref{tab:network} merinci alamat IP dan port setiap
\textit{node} dalam arsitektur. Setiap \textit{container} memiliki dua
port: port internal yang digunakan untuk komunikasi antar
\textit{container} dalam jaringan Docker, dan port \textit{host} yang
dipetakan untuk akses dari luar (digunakan selama pengembangan dan
pengujian). Aplikasi Laravel diekspos pada port host 8080 yang
dipetakan ke port 80 internal Nginx di dalam \textit{container}.
Pgpool-II membuka dua port: 5432 untuk koneksi database PostgreSQL dan
9898 untuk antarmuka administrasi PCP (\textit{Pgpool Control
Protocol}). Masing-masing PostgreSQL \textit{backend} diekspos pada
port host yang berbeda (5433, 5434, 5435) untuk memungkinkan koneksi
langsung saat debugging, meskipun dalam operasi normal aplikasi hanya
terhubung melalui Pgpool-II.

\begin{table}[H]
\centering
\caption{Spesifikasi Jaringan Tiap Node}
\label{tab:network}
\begin{tabularx}{\textwidth}{l l l l l}
\toprule
\textbf{Komponen} & \textbf{Container} & \textbf{IP Address} &
\textbf{Port (Host)} & \textbf{Port (Container)} \\
\midrule
Laravel App   & \texttt{bdt-app}     & 172.25.0.30 & 8080       & 80 \\
Pgpool-II     & \texttt{pgpool}      & 172.25.0.20 & 5436, 9898 & 5432, 9898 \\
PostgreSQL Master  & \texttt{pg-master}   & 172.25.0.10 & 5433       & 5432 \\
PostgreSQL Replica 1 & \texttt{pg-replica1} & 172.25.0.11 & 5434    & 5432 \\
PostgreSQL Replica 2 & \texttt{pg-replica2} & 172.25.0.12 & 5435    & 5432 \\
\bottomrule
\end{tabularx}
\end{table}

\subsection{Arah Komunikasi Antar Server}

Komunikasi antar \textit{node} mengikuti alur yang telah dirancang
untuk memastikan pemisahan beban kerja dan keamanan data. Berikut
adalah rincian arah komunikasi pada setiap jalur:

\begin{enumerate}
    \item \textbf{Client $\rightarrow$ Aplikasi (HTTP)}: Browser
    mengakses aplikasi Laravel melalui port \texttt{8080} pada
    \textit{host}. Permintaan HTTP ini kemudian diteruskan oleh Nginx
    ke aplikasi Laravel melalui PHP-FPM yang dikelola Supervisor.
    \item \textbf{Aplikasi $\rightarrow$ Pgpool-II (PostgreSQL Wire
    Protocol)}: Seluruh koneksi database dari aplikasi diarahkan ke
    Pgpool-II (172.25.0.20:5432), bukan langsung ke PostgreSQL Master.
    Aplikasi tidak perlu mengetahui topology database di belakang
    Pgpool-II karena middleware ini menangani seluruh logika routing
    query secara transparan.
    \item \textbf{Pgpool-II $\rightarrow$ PostgreSQL Master (WRITE)}:
    Query \texttt{INSERT}, \texttt{UPDATE}, \texttt{DELETE} diteruskan
    ke Master (172.25.0.10:5432). Pgpool-II mengenali jenis operasi
    SQL dan mengarahkannya ke \textit{backend} 0 (Master) yang
    memiliki flag \texttt{ALLOW\_TO\_FAILOVER}.
    \item \textbf{Pgpool-II $\rightarrow$ PostgreSQL Replica (READ)}:
    Query \texttt{SELECT} didistribusikan ke Replica 1
    (172.25.0.11:5432) dan Replica 2 (172.25.0.12:5432) berdasarkan
    \textit{weight} masing-masing. Dengan bobot yang sama (weight=1),
    kedua Replica mendapat porsi query yang seimbang.
    \item \textbf{Master $\rightarrow$ Replica (Streaming
    Replication)}: Data WAL (\textit{Write-Ahead Log}) dikirim dari
    Master ke kedua Replica secara \textit{streaming} untuk menjaga
    sinkronisasi data. Setiap perubahan yang terjadi di Master akan
    direplikasi ke Replica dalam waktu yang hampir \textit{real-time}.
    \item \textbf{Pgpool-II $\rightarrow$ Semua Backend (Health
    Check)}: Pgpool-II secara periodik memeriksa kesehatan setiap
    \textit{backend} untuk memastikan ketersediaan layanan. Jika suatu
    \textit{backend} gagal merespons setelah beberapa kali percobaan,
    \textit{backend} tersebut akan di-\textit{detach} dari pool.
\end{enumerate}

\section{Spesifikasi Lingkungan (Environment)}

Tabel~\ref{tab:environment} menjelaskan spesifikasi lingkungan
pengembangan dan \textit{deployment} yang digunakan dalam
implementasi. Seluruh komponen berjalan di atas Docker Engine pada
sistem operasi Linux (Ubuntu 22.04). PostgreSQL 16 dipilih karena
dukungan \textit{streaming replication} yang matang dan performa yang
teruji pada beban kerja tinggi. Pgpool-II versi terbaru
(\texttt{pgpool/pgpool:latest}) digunakan untuk memastikan
kompatibilitas penuh dengan PostgreSQL 16. Aplikasi Laravel berjalan
di atas PHP 8.2 dengan Nginx sebagai web server dan Supervisor sebagai
process manager untuk menjaga proses tetap berjalan.

\begin{table}[H]
\centering
\caption{Spesifikasi Lingkungan Pengembangan dan Deployment}
\label{tab:environment}
\begin{tabularx}{\textwidth}{l X}
\toprule
\textbf{Komponen} & \textbf{Spesifikasi} \\
\midrule
Operating System (Host)  & Ubuntu 22.04 / Linux (via Docker) \\
Database                 & PostgreSQL 16 \\
Middleware               & Pgpool-II 4.5
                           (\texttt{pgpool/pgpool:latest}) \\
Deployment               & Docker Compose v2+ (5 container) \\
Bahasa Pemrograman       & PHP 8.2 \\
Framework                & Laravel 10 \\
Web Server               & Nginx (Alpine Linux) \\
Process Manager          & Supervisor \\
Network                  & Docker Bridge
                           (\texttt{bdt-network}, 172.25.0.0/24) \\
Replikasi                & PostgreSQL Streaming Replication
                           (Async) \\
Volume Persistence       & Docker Named Volumes
                           (\texttt{pg\_master\_data},
                            \texttt{pg\_replica1\_data},
                            \texttt{pg\_replica2\_data}) \\
\bottomrule
\end{tabularx}
\end{table}

\subsection{Struktur File Konfigurasi}

Struktur direktori proyek diorganisasikan untuk memisahkan konfigurasi
infrastruktur database dari kode aplikasi. Direktori
\texttt{docker/} berisi seluruh konfigurasi yang spesifik untuk
masing-masing layanan database, sementara direktori
\texttt{PWEB-B1/} menyimpan kode aplikasi Laravel beserta konfigurasi
\textit{deployment}-nya. Berikut adalah struktur direktori proyek yang
relevan dengan konfigurasi basis data terdistribusi:

{\footnotesize
\begin{verbatim}
BDT/
+-- docker-compose.yml                  # Orchestration semua service
+-- docker/
|   +-- pg-master/
|   |   +-- postgresql.conf             # Konfigurasi PostgreSQL Master
|   |   +-- pg_hba.conf                 # Host-Based Authentication
|   |   +-- init.sql                    # Inisialisasi user & role
|   +-- pg-replica/
|   |   +-- postgresql.conf             # Konfigurasi PostgreSQL Replica
|   |   +-- entrypoint.sh               # Script pg_basebackup otomatis
|   +-- pgpool/
|       +-- pgpool.conf                 # Konfigurasi Pgpool-II
|       +-- pool_hba.conf               # Autentikasi client Pgpool
|       +-- pcp.conf                    # PCP admin user
+-- PWEB-B1/
    +-- Dockerfile                      # Build image Laravel
    +-- .env                            # Environment (DB_HOST=pgpool)
    +-- .env.bdt                        # Environment Docker
    +-- docker/
    |   +-- nginx/default.conf          # Konfigurasi Nginx
    |   +-- supervisor/supervisord.conf # Process manager
    +-- config/database.php             # Konfigurasi koneksi database
\end{verbatim}
}
